\documentclass[11pt]{article}
\usepackage[margin=1in]{geometry}
\usepackage{amsmath,amssymb}
\usepackage{graphicx}
\usepackage{booktabs}
\usepackage{xcolor}
\usepackage{listings}
\usepackage{hyperref}
\usepackage{siunitx}
\usepackage{caption}

\definecolor{codebg}{RGB}{247,247,247}
\definecolor{codekw}{RGB}{0,90,160}
\definecolor{codestr}{RGB}{140,20,20}
\definecolor{codecmt}{RGB}{90,90,90}

\lstdefinestyle{py}{
  backgroundcolor=\color{codebg},
  basicstyle=\ttfamily\footnotesize,
  keywordstyle=\color{codekw}\bfseries,
  commentstyle=\color{codecmt}\itshape,
  stringstyle=\color{codestr},
  numbers=left,
  numberstyle=\tiny\color{gray},
  stepnumber=1,
  numbersep=8pt,
  breaklines=true,
  showstringspaces=false,
  frame=single,
  rulecolor=\color{gray!40},
  language=Python,
  tabsize=4,
  captionpos=b
}

\title{\vspace{-1.5cm}ME698V: Physics-Based Modeling of Li-ion Batteries \\[4pt]
\large Mid-Term Project \\[2pt]
\large Development of a Single-Particle Model (SPM) for \\
Carbon $|$ LiPF\textsubscript{6} Electrolyte $|$ NMC111 Li-ion Cells}
\author{}
\date{}

\begin{document}
\maketitle
\vspace{-1cm}

\section*{Project Statement}
A commercial electric vehicle (EV) manufacturer requires a mathematical model for the Carbon $|$ LiPF\textsubscript{6}
Electrolyte $|$ NMC111 cell used in their EV powertrain. A single-particle model (SPM) is developed in three stages:
(i) the Carbon $|$ LiPF\textsubscript{6} $|$ Li half-cell is modeled following Verbrugge and Koch's thermodynamic
interaction formalism; (ii) the NMC111 $|$ LiPF\textsubscript{6} $|$ Li half-cell is modeled with an analogous
formalism and its interaction parameters fit to the reported voltage-capacity data; (iii) the two half-cell models
are combined into a full-cell response, first at equilibrium and then including Butler--Volmer kinetics at the
particle surface.

\section{Carbon Anode Half-Cell}

\subsection{Chemical potential and governing equation}
Following Verbrugge and Koch, the chemical potential of intercalated lithium in the carbon host is
\begin{equation}
\mu_{\mathrm{Li}} = \mu_{\mathrm{Li}}^{0} + RT\ln\!\left(\frac{c}{1-c}\right) + \sum_{k=2}^{7} k\,\Omega_k^{a}\,c^{\,k-1},
\qquad 0\le c \le 1,
\end{equation}
where $c$ is the fractional occupancy and $\Omega_k^{a}$ are the anode interaction coefficients. The flux is driven
by the chemical-potential gradient,
\begin{equation}
N = -\frac{Dc}{RT}\nabla \mu_{\mathrm{Li}} = -D\!\left(c\,\frac{d\mu_{\mathrm{Li}}}{dc}\right)\nabla c .
\end{equation}
Differentiating,
\begin{equation}
c\,\frac{d\mu_{\mathrm{Li}}}{dc} = \frac{RT}{1-c} + \sum_{k=2}^{7} k(k-1)\,\Omega_k^{a}\,c^{\,k-1},
\end{equation}
so that $N = -D_{\mathrm{eff}}(c)\,\nabla c$ with the thermodynamic-factor-corrected diffusivity
\begin{equation}
D_{\mathrm{eff}}(c) = D\left[\frac{1}{1-c} + \frac{1}{RT}\sum_{k=2}^{7} k(k-1)\,\Omega_k^{a}\,c^{\,k-1}\right].
\label{eq:Deff_anode}
\end{equation}
For the single particle (spherical geometry, radial transport only), the material balance is
\begin{equation}
\frac{\partial c}{\partial t} = -\nabla\cdot N = \frac{1}{r^2}\frac{\partial}{\partial r}
\left(r^2 D_{\mathrm{eff}}(c)\,\frac{\partial c}{\partial r}\right).
\label{eq:pde_anode}
\end{equation}

\subsection{Boundary and initial conditions}
Symmetry at the particle center and continuity of the intercalation flux with the imposed current density
$i_{\mathrm{app}}$ at the surface ($i_{\mathrm{app}}>0$ denotes Li insertion / lithiation) give
\begin{align}
\left.\frac{\partial c}{\partial r}\right|_{r=0} &= 0, \\
\left.\frac{\partial c}{\partial r}\right|_{r=R} &= \frac{i_{\mathrm{app}}}{F\,D_{\mathrm{eff}}(c_s)}, \qquad c_s \equiv c(R,t),
\label{eq:bc_anode}
\end{align}
with initial condition $c(r,0) = c_0 = 5.0\times10^{-4}$ (uniform).

\subsection{Cell voltage}
With the exchange current density taken large compared to $i_{\mathrm{app}}$ ($\eta \approx 0$), the cell voltage
relative to a Li metal reference is
\begin{equation}
V = U^{s} + \frac{RT}{F}\ln\!\left(\frac{1-c_s}{c_s}\right) - \frac{1}{F}\sum_{k=2}^{7} k\,\Omega_k^{a}\,c_s^{\,k-1}.
\label{eq:V_anode}
\end{equation}

\subsection{Parameters and implementation}
Table~\ref{tab:anode_params} lists the parameters used, taken from Tables I and II of Verbrugge and Koch (1996).
Equations~\eqref{eq:Deff_anode}--\eqref{eq:V_anode} were implemented in PyBaMM as a custom \texttt{BaseModel} on a
spherical particle domain, discretized with the finite-volume method (50 nodes) and integrated with the default
solver for $i_{\mathrm{app}} = +1.0~\si{mA/cm^2}$.

\begin{table}[h]
\centering
\caption{Carbon anode parameters (Verbrugge and Koch, 1996).}
\label{tab:anode_params}
\begin{tabular}{lll}
\toprule
Parameter & Symbol & Value \\
\midrule
Particle radius & $R$ & \SI{3.5}{\micro\metre} \\
Site concentration & $c_{\max}$ & \SI{18000}{mol/m^3} \\
Guest diffusivity & $D$ & $1\times10^{-14}~\si{m^2/s}$ \\
Standard potential & $U^{s}$ & \SI{0.8170}{V} \\
$\Omega_2^{a}/F,\ldots,\Omega_7^{a}/F$ & --- & $0.9926,\ 0.8981,\ -5.630,\ 8.595,\ -5.784,\ 1.468$~V \\
\bottomrule
\end{tabular}
\end{table}

\begin{figure}[h]
\centering
\includegraphics[width=0.72\linewidth]{Problem1_updated.png}
\caption{Simulated and experimental (Table 1) voltage vs. average occupancy for the carbon half-cell,
$i_{\mathrm{app}} = 1.0~\si{mA/cm^2}$.}
\label{fig:p1}
\end{figure}

The simulated charging curve reproduces the experimental data of Fig.~1(a) closely across the full occupancy range
(Fig.~\ref{fig:p1}), confirming that the thermodynamic-factor-corrected diffusivity of Eq.~\eqref{eq:Deff_anode}
combined with the equilibrium voltage relation of Eq.~\eqref{eq:V_anode} captures the strongly non-ideal
intercalation behavior reported by Verbrugge and Koch.
\clearpage
\section{NMC111 Cathode Half-Cell}

\subsection{Chemical potential and governing equation}
The chemical potential of lithium in NMC111 is written with the interaction sum evaluated in the vacancy fraction,
\begin{equation}
\mu_{\mathrm{Li}} = \mu_{\mathrm{Li}}^{0} + RT\ln\!\left(\frac{c}{1-c}\right) - \sum_{k=2}^{7} k\,\Omega_k^{c}\,(1-c)^{\,k-1},
\end{equation}
where $c$ is again the fractional occupancy of Li in the host. Proceeding exactly as in Section 1,
\begin{equation}
c\,\frac{d\mu_{\mathrm{Li}}}{dc} = \frac{RT}{1-c} + \sum_{k=2}^{7} k(k-1)\,\Omega_k^{c}\;c\,(1-c)^{\,k-2},
\end{equation}
so the effective diffusivity for the cathode particle is
\begin{equation}
D_{\mathrm{eff}}(c) = D\left[\frac{1}{1-c} + \frac{1}{RT}\sum_{k=2}^{7} k(k-1)\,\Omega_k^{c}\;c\,(1-c)^{\,k-2}\right].
\label{eq:Deff_cathode}
\end{equation}
Note the extra factor of $c$ multiplying $(1-c)^{k-2}$: because the interaction sum in $\mu_{\mathrm{Li}}$ is a
function of $(1-c)$ rather than $c$, the chain rule introduces this asymmetry relative to the anode's
Eq.~\eqref{eq:Deff_anode}, and the two expressions are \emph{not} related by the naive substitution $c\to(1-c)$.
The governing PDE and symmetry condition at $r=0$ retain the same form as Eqs.~\eqref{eq:pde_anode} and the first
line of \eqref{eq:bc_anode}. The surface flux condition, initial condition, and voltage expression are
\begin{align}
\left.\frac{\partial c}{\partial r}\right|_{r=R} &= \frac{i_{\mathrm{app}}}{F\,D_{\mathrm{eff}}(c_s)}, \qquad
c(r,0) = 0.99, \\[4pt]
V &= V_0 + \frac{RT}{F}\ln\!\left(\frac{1-c_s}{c_s}\right) + \frac{1}{F}\sum_{k=2}^{7} k\,\Omega_k^{c}\,(1-c_s)^{\,k-1}.
\label{eq:V_cathode}
\end{align}
With $i_{\mathrm{app}} = -1.0~\si{mA/cm^2}$ (Li\textsuperscript{+} leaving the cathode), $c_s$ decreases from 0.99 and
the specific capacity referenced to the NMC111 mass is
\begin{equation}
Q_{\mathrm{sp}}(t) = \frac{F\,c_{\max}}{3600\,\rho_{\mathrm{NMC}}}\big(c_0 - c_s(t)\big)\ \ [\si{mAh/g}],
\qquad
\frac{F\,c_{\max}}{3600\,\rho_{\mathrm{NMC}}} = \SI{609.7}{mAh/g} \equiv Q_{\max}.
\label{eq:Qsp}
\end{equation}

\subsection{Parameter fitting}
$V_0$ and $\Omega_k^c$ ($k=2,\ldots,7$) were obtained by a constrained nonlinear least-squares fit
(\texttt{scipy.optimize.curve\_fit}) of the equilibrium ($\eta = 0$) voltage--capacity relation, Eqs.~\eqref{eq:V_cathode}
and \eqref{eq:Qsp}, to the twenty data points of Table 2, with bounds fixed to the ranges quoted in the project
statement:
\[
V_0\in[3.0,4.0],\quad \Omega_2^c/F\in[-0.15,0],\quad \Omega_3^c/F\in[1.5,2.0],\quad \Omega_4^c/F\in[-3.2,-2.5],
\]
\[
\Omega_5^c/F\in[10.0,14.0],\quad \Omega_6^c/F\in[11,15],\quad \Omega_7^c/F\in[-4.0,-2.5].
\]
The fit converged to the values in Table~\ref{tab:fit}, with an RMS residual of \SI{22}{mV} over the fitted range;
the same optimum was recovered from several different starting points inside the bounds.

\begin{table}[h]
\centering
\caption{Fitted NMC111 interaction parameters.}
\label{tab:fit}
\begin{tabular}{lc}
\toprule
Parameter & Fitted value \\
\midrule
$V_0$ & \SI{3.824}{V} \\
$\Omega_2^c/F$ & $-0.000$~V \\
$\Omega_3^c/F$ & $2.000$~V \\
$\Omega_4^c/F$ & $-2.515$~V \\
$\Omega_5^c/F$ & $10.000$~V \\
$\Omega_6^c/F$ & $11.000$~V \\
$\Omega_7^c/F$ & $-4.000$~V \\
\midrule
RMSE & \SI{22}{mV} \\
\bottomrule
\end{tabular}
\end{table}

\subsection{Transient SPM simulation}
With the fitted parameters and the effective diffusivity of Eq.~\eqref{eq:Deff_cathode}, the transient PDE was
solved with $D = 7\times10^{-15}~\si{m^2/s}$, $R_p = \SI{10}{\micro\metre}$, $c_{\max} = \SI{48000}{mol/m^3}$,
$\rho_{\mathrm{NMC}} = \SI{2.11}{g/cm^3}$, and $i_{\mathrm{app}} = -1.0~\si{mA/cm^2}$.

\begin{figure}[h]
\centering
\includegraphics[width=0.72\linewidth]{Problem2_updated.png}
\caption{Simulated (fitted-parameter SPM) and experimental (Table 2) voltage vs. specific capacity for the
NMC111 half-cell.}
\label{fig:p2}
\end{figure}
\clearpage
\section{Full-Cell Voltage vs. Specific Capacity}

\subsection{Combining the half-cells}
The full cell is obtained by running the anode model of Section 1 ($i_{\mathrm{app}} = +1.0~\si{mA/cm^2}$) and the
cathode model of Section 2 ($i_{\mathrm{app}} = -1.0~\si{mA/cm^2}$) \emph{concurrently over the same physical time},
consistent with the stated full-cell current $i_{\mathrm{app}} = 1.0~\si{mA/cm^2}$ (positive from cathode to anode).
Because the same current density is applied to both electrodes for the same elapsed time, charge is conserved
between them automatically -- no independent electrode-mass or loading ratio is required beyond the single-particle
parameters already given in Problems 1 and 2. The full-cell voltage is then
\begin{equation}
V_{\mathrm{full}}(t) = V_{\mathrm{cathode}}(t) - V_{\mathrm{anode}}(t),
\end{equation}
reported against the cathode's own specific-capacity trace $Q_{\mathrm{sp}}(t)$ from Eq.~\eqref{eq:Qsp}. The curve is
physically meaningful up to the point at which either electrode is exhausted; here the anode (the smaller-capacity
electrode, given its particle size and site concentration) is the limiting electrode, and the trace is shown up to
$\theta_{\mathrm{anode}}\to 1$.

\begin{figure}[h]
\centering
\includegraphics[width=0.72\linewidth]{Problem3_updated.png}
\caption{Full-cell voltage vs. specific capacity (NMC111 basis), obtained by combining the anode and cathode SPMs of
Problems 1 and 2 at $i_{\mathrm{app}} = 1.0~\si{mA/cm^2}$.}
\label{fig:p3}
\end{figure}

The full-cell curve (Fig.~\ref{fig:p3}) rises smoothly from about \SI{2.9}{V} at low capacity to over \SI{4.5}{V} as
the anode approaches saturation, with the characteristic steep upturn near \SIrange{150}{175}{mAh/g} that reflects
the anode voltage collapsing toward zero as $\theta_{\mathrm{anode}}\to1$ -- the expected signature of the
capacity-limiting electrode in a full cell.
\clearpage
\section{Cyclic Full-Cell Voltage with Finite Exchange Current Density}

\subsection{Butler--Volmer overpotential}
For symmetry factor $\beta = 1/2$, the Butler--Volmer relation $i = i_{\mathrm{ex}}\big[e^{(1-\beta)f\eta} - e^{-\beta f\eta}\big]$
inverts to
\begin{equation}
\eta = \frac{2RT}{F}\,\operatorname{asinh}\!\left(\frac{i_{\mathrm{app}}}{2\,i_{\mathrm{ex}}}\right).
\label{eq:eta}
\end{equation}
Both electrode/electrolyte interfaces are assigned the same $i_{\mathrm{ex}}$. On \emph{discharge}
($\mathrm{Li^+}$ moving cathode $\to$ anode reversed relative to Problem 3, i.e. the cell current sourced externally),
each electrode reaction is displaced below its equilibrium potential in the direction that opposes the current, so
the cell voltage is reduced relative to open circuit by twice the single-electrode overpotential; on \emph{charge}
the current reverses and the same overpotential instead \emph{adds} to the open-circuit voltage:
\begin{equation}
V_{\mathrm{discharge}}(Q) = V_{\mathrm{OCV}}(Q) - 2\eta, \qquad
V_{\mathrm{charge}}(Q) = V_{\mathrm{OCV}}(Q) + 2\eta,
\label{eq:cycle}
\end{equation}
where $V_{\mathrm{OCV}}(Q)$ is the equilibrium full-cell curve of Fig.~\ref{fig:p3} (Section 3, $\eta=0$). One full
cycle is constructed by evaluating Eq.~\eqref{eq:cycle} along the discharge leg from $Q=0$ up to the capacity limit,
then back along the reversed state-of-charge path for the charge leg, plotted against cumulative throughput capacity.

\subsection{Results}
Figure~\ref{fig:p4} shows the resulting cycle for $i_{\mathrm{ex}} \in \{1, 10, 10^2, 10^3, 10^4\}~\si{mA/cm^2}$ at
$|i_{\mathrm{app}}| = 1.0~\si{mA/cm^2}$. As Eq.~\eqref{eq:eta} predicts, $\eta \to 0$ once $i_{\mathrm{ex}} \gg
i_{\mathrm{app}}$, so the $i_{\mathrm{ex}}\ge 10~\si{mA/cm^2}$ curves collapse onto the equilibrium curve and are
indistinguishable from one another; only $i_{\mathrm{ex}} = 1~\si{mA/cm^2}$ (comparable in magnitude to
$i_{\mathrm{app}}$) shows a visibly open hysteresis loop, with the charge leg running above and the discharge leg
below the equilibrium curve by $2\eta \approx \SI{50}{mV}$ near mid-capacity.

\begin{figure}[h]
\centering
\includegraphics[width=0.85\linewidth]{Problem4_updated.png}
\caption{Cyclic full-cell voltage vs. throughput capacity for five exchange current densities,
$|i_{\mathrm{app}}| = 1.0~\si{mA/cm^2}$.}
\label{fig:p4}
\end{figure}
\clearpage
\section*{Summary}
A single-particle model was developed for the Carbon $|$ LiPF\textsubscript{6} $|$ NMC111 cell. The carbon anode
half-cell (Section 1) reproduces the Verbrugge--Koch interaction-model voltage--occupancy curve directly from
published parameters. The NMC111 half-cell (Section 2) uses an analogous thermodynamic-factor-corrected diffusivity,
with its seven open interaction parameters obtained by constrained least-squares fitting to the reported
voltage--capacity data (RMSE \SI{22}{mV}). Combining the two half-cell SPMs under a shared applied current
(Section 3) yields a full-cell voltage--capacity curve that correctly reflects the anode as the capacity-limiting
electrode. Finally, incorporating Butler--Volmer kinetics at finite exchange current density (Section 4) reproduces
the expected narrowing of the charge/discharge hysteresis loop as $i_{\mathrm{ex}}$ increases, collapsing to the
equilibrium curve once the exchange current density greatly exceeds the applied current density.

\end{document}
